
Below is an updated executive architecture summary that reframes the system as an Agentic Experimenter for Microscopic Traffic Simulation.
This version aligns with the full discussion we had, but now explicitly incorporates the agent loop, tool usage, DSL layer, invariants, and simulation workflow.

This format is written at a level suitable for:

\begin{itemize}
\item dissertation planning


\item research proposals


\item system architecture documentation


\end{itemize}
\chapter*{Executive Architecture Summary}
\section*{Agentic Experimenter for Microscopic Traffic Simulation}
\subsection*{Target Application}
Wildfire evacuation and contraflow evaluation on the I-10 corridor during the Palisades Fire scenario using a data-calibrated microscopic traffic simulator.

The system enables an LLM-driven agent to autonomously design and execute simulation experiments while preserving physical validity and simulation correctness through structured constraints and validation layers.

The architecture integrates:

\begin{itemize}
\item real traffic data


\item knowledge graphs


\item LLM reasoning


\item DSL-based scenario generation


\item simulation execution


\item automated evaluation


\end{itemize}
to support agentic scientific experimentation on transportation systems.

\chapter*{1. Research Objective}
Develop an Agentic Experimenter system capable of autonomously exploring evacuation strategies in a microscopic traffic simulation environment.

The agent should be able to:

• design simulation scenarios
• modify evacuation strategies
• evaluate contraflow policies
• analyze evacuation performance
• iteratively refine experiments

The system acts as a simulation research assistant capable of running structured traffic experiments.

\chapter*{2. Foundational Simulation Engine}
The simulation core is implemented using the ScalaTion microscopic traffic simulation framework.

The engine models:

• vehicles
• lanes
• merges
• event scheduling
• car-following behavior
• lane-changing dynamics
• network topology

The simulation engine is trusted infrastructure and is never modified by the LLM.

This separation ensures that simulation physics and causality remain protected.

\chapter*{3. Real-World Data Sources}
The system uses empirical traffic datasets to calibrate and reconstruct realistic traffic dynamics.

\subsection*{Primary Dataset}
Caltrans PeMS traffic data

Provides:

• lane-level flow
• speed
• occupancy
• sensor metadata

These data support:

• baseline calibration
• wildfire-day congestion reconstruction
• demand surge estimation.

\chapter*{4. Transportation Network Representation}
Road network structure is derived from:

OpenStreetMap (OSM).

The network may be stored in a Neo4j knowledge graph, representing:

Nodes:

• road segments
• lanes
• sensors
• ramps

Edges:

• connectivity
• lane adjacency
• merges
• corridor topology.

This representation enables structured reasoning about network structure.

\chapter*{5. Graph-Based Knowledge Retrieval}
The agent can access the knowledge graph through GraphRAG (Graph Retrieval Augmented Generation).

GraphRAG allows the LLM to retrieve contextual information such as:

• corridor geometry
• lane counts
• sensor placement
• ramp connections.

This ensures that generated simulation scenarios remain consistent with network structure.

\chapter*{6. DSL-Based Scenario Generation}
To prevent unsafe code generation, the LLM does not generate simulation code directly.

Instead, it generates a structured Scala DSL describing simulation scenarios.

Example scenario specification:

scenario "Wildfire_I10" {
  corridor "I10_EB"
  surge at "10:30" scale 1.8
  smoke moderate
  contraflow +1 preserveInbound 1
}
The DSL defines:

• demand surge conditions
• smoke-induced behavioral degradation
• contraflow lane allocation
• corridor selection.

The DSL is interpreted by the simulation engine.

\chapter*{7. Scenario Validation Layer}
Before simulation execution, DSL scenarios undergo validation checks.

These checks ensure:

• valid lane configurations
• preservation of inbound emergency lanes
• corridor existence
• valid parameter ranges
• structural consistency.

Invalid scenarios are rejected before simulation begins.

\chapter*{8. Agentic Experiment Loop}
The core of the architecture is an agentic experimentation loop.

The agent repeatedly designs and evaluates simulation scenarios.

define research objective
↓
generate simulation scenario (DSL)
↓
validate scenario
↓
run simulation
↓
collect metrics
↓
analyze results
↓
propose new experiment
↓
repeat
This loop allows the agent to search the experiment space autonomously.

\chapter*{9. Tool-Based Agent Architecture}
The agent interacts with the system through structured tools.

\subsection*{Scenario Generation Tool}
Produces DSL scenario configurations.

\subsection*{Network Query Tool}
Retrieves topology information from Neo4j.

\subsection*{Simulation Runner Tool}
Executes the traffic simulation.

\subsection*{Validator Tool}
Checks scenario consistency.

\subsection*{Metrics Tool}
Computes evacuation performance metrics.

These tools enable structured reasoning and experimentation.

\chapter*{10. Runtime Semantic Safety Layer}
Simulation correctness cannot be guaranteed through static analysis alone.

The system therefore enforces runtime invariants to ensure physical plausibility.

Examples include:

\subsection*{Vehicle collision prevention}
Vehicles must maintain minimum spacing.

\subsection*{Speed validity}
Speeds must remain finite and non-negative.

\subsection*{Flow conservation}
Vehicle counts must remain consistent.

\subsection*{Merge ordering}
Vehicle ordering must remain physically valid.

\subsection*{Event queue integrity}
Simulation events must remain time-ordered.

These checks prevent semantic corruption of simulation results.

\chapter*{11. Static Code Analysis Support}
If DSL-generated code or extensions are produced, static analysis tools may be used to detect unsafe patterns.

Scala tools include:

• Scalafix
• Scapegoat
• WartRemover.

These tools detect:

• unsafe constructs
• null usage
• problematic code patterns.

However, they complement but do not replace runtime simulation invariants.

\chapter*{12. Experimental Workflow}
The agent conducts experiments related to wildfire evacuation dynamics.

The experiment phases include:

\subsection*{Phase 1 — Baseline Calibration}
Reconstruct normal traffic dynamics.

\subsection*{Phase 2 — Fire-Day Reconstruction}
Model demand surge from wildfire evacuation.

\subsection*{Phase 3 — Smoke Impact Modeling}
Simulate behavioral degradation due to smoke.

\subsection*{Phase 4 — Contraflow Evaluation}
Test counterfactual evacuation strategies.

\chapter*{13. Evaluation Metrics}
Simulation performance is evaluated using traffic engineering metrics.

These include:

• throughput (vehicles/hour)
• mean speed
• congestion clearance time
• shockwave propagation speed
• resilience index.

These metrics quantify evacuation performance.

\chapter*{14. Architecture Overview}
The complete system architecture can be summarized as:

Real Traffic Data (PeMS)
        │
        ▼
Knowledge Graph (Neo4j + OSM)
        │
        ▼
GraphRAG Retrieval
        │
        ▼
LLM Agent
        │
        ▼
Scenario DSL Generation
        │
        ▼
Scenario Validation
        │
        ▼
Microscopic Traffic Simulation
        │
        ▼
Runtime Invariant Checks
        │
        ▼
Performance Metrics
        │
        ▼
Agent Analysis and Next Experiment
\chapter*{15. Advantages of the Agentic Approach}
The proposed system provides several benefits.

\subsection*{Autonomous Experimentation}
The agent can explore evacuation strategies without manual configuration.

\subsection*{Safety and Correctness}
DSL constraints and invariants protect simulation integrity.

\subsection*{Data-Driven Modeling}
Real traffic data ensures realistic scenarios.

\subsection*{Scalable Experiment Design}
Thousands of evacuation scenarios can be evaluated.

\subsection*{Scientific Discovery}
The system can identify evacuation strategies that improve resilience.

\chapter*{16. Long-Term Extensions}
Future developments may include:

• multi-corridor evacuation modeling
• adaptive contraflow optimization
• integration with reinforcement learning
• real-time digital twin systems
• multi-agent experimentation.

\chapter*{Final Positioning Statement}
This architecture introduces an LLM-driven Agentic Experimenter capable of autonomously generating and evaluating wildfire evacuation strategies in a microscopic traffic simulation environment.

By combining structured DSL scenario generation, graph-based knowledge retrieval, simulation execution, and runtime invariant enforcement, the system enables safe and scalable exploration of evacuation policies while preserving simulation validity.

If you'd like, I can also show you the three strongest papers you could write from this architecture, because you actually have multiple publishable angles hiding inside this design.

"""""
Paper angles:

You actually have at least three distinct, publishable papers hidden in the architecture we just built. Each paper focuses on a different research contribution, even though they share the same system.

I’ll show them as Paper A, B, C, along with venue fit and the core contribution.

\chapter*{Paper A — LLM Agent for Traffic Simulation Experiment Design}
\section*{(The Agentic Experimenter Paper)}
\subsection*{Core Idea}
Demonstrate that an LLM agent can autonomously design and execute traffic simulation experiments to explore evacuation strategies.

The novelty is not the traffic simulation itself — it is the agentic experimentation loop.

\section*{Problem}
Designing evacuation experiments manually is slow and biased.

Researchers typically test only a small number of scenarios.

Example:

contraflow +1 lane
contraflow +2 lanes
no contraflow
But the scenario space is huge.

\section*{Proposed System}
Your Agentic Experimenter:

Goal:
Minimize evacuation clearance time
Agent loop:

generate scenario
↓
run simulation
↓
evaluate metrics
↓
propose new scenario
↓
repeat
The agent explores combinations like:

+1 contraflow
+2 contraflow
+1 contraflow + moderate smoke
+2 contraflow + severe smoke
demand surge × 1.5
demand surge × 2
\section*{Contribution}
You show that:

LLM agents can autonomously explore evacuation strategies in traffic simulation environments.

\section*{Evaluation}
Compare:

Human-designed experiments
vs
Agent-generated experiments
Measure:

\begin{itemize}
\item best evacuation strategy discovered


\item search efficiency


\item coverage of scenario space


\end{itemize}
\section*{Good Venues}
• ANNSIM
• Winter Simulation Conference (WSC)
• IEEE ITSC

\chapter*{Paper B — Safe LLM Integration into Simulation Systems}
\section*{(The DSL + Invariant Paper)}
\subsection*{Core Idea}
LLMs generating simulation code is dangerous because simulations can compile but produce invalid physics.

Your solution:

LLM generates a DSL instead of raw code, combined with runtime invariants.

\section*{Problem}
Simulation correctness is semantic, not syntactic.

Even perfect code can produce:

\begin{itemize}
\item vehicle overlap


\item NaN speeds


\item flow conservation violations


\item broken event queues


\end{itemize}
Static analysis cannot catch these.

\section*{Proposed Architecture}
LLM
↓
DSL scenario
↓
validator
↓
simulation engine
↓
runtime invariants
↓
metrics
Safety layers:

\begin{enumerate}
\item DSL constraints


\item scenario validation


\item runtime invariant checks


\end{enumerate}
\section*{Contribution}
You propose a safe architecture for integrating LLMs into simulation systems.

Key idea:

LLMs should generate structured scenario specifications, not simulation code.

\section*{Evaluation}
Demonstrate:

Without DSL:

LLM generated invalid simulations
With DSL:

invalid scenarios rejected
\section*{Good Venues}
• WSC (Winter Simulation Conference)
• ACM SIGSIM PADS
• Simulation Modelling Practice and Theory

\chapter*{Paper C — Wildfire Evacuation Contraflow Analysis}
\section*{(The Domain Study)}
This is your transportation paper.

It uses the system but focuses on traffic science.

\section*{Research Question}
Does contraflow improve evacuation resilience
during wildfire-induced congestion?
Study:

I-10 corridor during Palisades Fire scenario

\section*{Experimental Design}
Scenarios:

A: no contraflow
B: +1 reversed lane
C: +2 reversed lanes
D: contraflow + smoke degradation
Metrics:

\begin{itemize}
\item evacuation throughput


\item congestion clearance time


\item shockwave speed


\item resilience index


\end{itemize}
\section*{Contribution}
Provide data-calibrated evidence on contraflow effectiveness in wildfire evacuation.

\section*{Good Venues}
• Winter Simulation Conference
• Transportation Research Part C
• IEEE ITSC

\chapter*{The Nice Thing}
All three papers share the same system.

Agent system
↓
DSL architecture
↓
contraflow experiments
You just emphasize different aspects.

\chapter*{Strategic Publication Order}
Best order:

\subsection*{1️⃣ Contraflow Study}
(WSC)

Fastest publishable.

\subsection*{2️⃣ DSL Safety Architecture}
(Simulation conference)

More systems-oriented.

\subsection*{3️⃣ Agentic Experimenter}
(AI + simulation)

Most novel but hardest.

\chapter*{What Makes Your Work Interesting}
You are combining:

traffic simulation
+
LLM agents
+
knowledge graphs
+
DSL safety
Few people are working at this intersection.

\chapter*{Important Advice}
Do not try to publish everything in one paper.

Each paper should have one clean idea.



"""
All the above can become: NOT part of the current proposal.
Traffic Data (PeMS)
│
▼
Network Graph (OSM + Neo4j)
│
▼
Causal World Model (SCM)
│
▼
Agent Reasoning
│
▼
Scenario DSL
│
▼
Simulation Engine (ScalaTion)
│
▼
Runtime Invariants
│
▼
Metrics + Observations
│
▼
World Model Update